\subsection{A maximum-rank comparison criterion}
\label{sec:nef-maximum-rank-comparison}
\index{Nef degeneration!maximum-rank comparison}

The constant-rank hypothesis in \cref{lem:nef-logarithmic-rank-comparison} can be removed. For a closed semipositive form $B$, the intersection polynomial $\int_X(B+\epsilon\Omega)^n$ has a leading term of order $\epsilon^{n-R}$, where $R$ is the largest rank that $B$ attains anywhere on $X$. Combined with the comparison of quotients, this gives the order $n-R$ for every background that is comparable to $B$. The criterion applies to backgrounds that vanish on an open set, and to positive variable multiples of arbitrary closed semipositive generators.

\Needspace{12\baselineskip}
\begin{proposition}\label{prop:nef-maximum-rank-comparison}
Let $X$ be a compact K\"ahler manifold of dimension $n$, let $B$ be a smooth closed semipositive $(1,1)$-form, and put $R=\max_X\operatorname{rank}B$. Suppose a smooth nef background $\theta$ satisfies
\[
 \theta\leq CB,\qquad \theta+\ddc\varphi\geq cB
\]
as currents, for constants $c,C>0$ and a real function $\varphi\in L^1(X)$. For every smooth positive probability form $\mu$ and Hermitian form $h$, there are constants $L,U>0$ such that
\[
 L\epsilon^{n-R}\leq\mathcal V_\mu(\theta+\epsilon h)
                       \leq U\epsilon^{n-R},\qquad 0<\epsilon\leq1.
\]
Thus $\theta$ has common order $n-R$, without any constant-rank hypothesis.
\end{proposition}
\begin{proof}
Fix a K\"ahler form $\Omega$. The first part of \cref{lem:nef-logarithmic-rank-comparison}, which does not use the rank, shows that $\lambda_\mu(\theta+\epsilon h)$ and $\lambda_\mu(B+\epsilon\Omega)$ differ by bounded positive factors. Since $B^{R+1}=0$, the closed nef formula following \cref{thm:gaud-volume-nef-limit} gives
\[
 \mathcal V_\mu(B+\epsilon\Omega)
 =\sum_{k=0}^{R}\binom nk\epsilon^{n-k}
                         \int_X B^k\wedge\Omega^{n-k}.
\]
All coefficients are nonnegative. If $R>0$, then at a point where $B$ has rank $R$, some $R\times R$ minor of its coefficient matrix is positive and stays positive nearby, so $B^R\wedge\Omega^{n-R}$ has positive integral. If $R=0$, this integral is the K\"ahler volume. Hence the coefficient of $\epsilon^{n-R}$ is positive. For $0<\epsilon\leq1$, every other power of $\epsilon$ that occurs is at most $\epsilon^{n-R}$, so the polynomial is bounded below by its leading term and above by the sum of its coefficients times $\epsilon^{n-R}$. Raising the comparison of quotients to the $n$th power gives the asserted bounds, and taking logarithms gives the order.
\end{proof}

\Needspace{12\baselineskip}
\begin{corollary}\label{cor:nef-variable-semipositive-generators}
On a compact K\"ahler manifold $X$ of dimension $n$, fix finitely many smooth closed semipositive $(1,1)$-forms $B_j$ and smooth strictly positive functions $f_j$. For $t_j\geq0$, put
\[
 \theta=\sum_jt_jf_jB_j,\qquad
 B=\sum_{j:t_j>0}B_j,\qquad R=\max_X\operatorname{rank}B.
\]
The background $\theta$ is nef, has common order $n-R$, and satisfies the two-sided power bounds of \cref{prop:nef-maximum-rank-comparison}. For every modification $\pi:Y\to X$ with $Y$ a compact K\"ahler manifold, the pullback $\pi^*\theta$ has the same order and power bounds with the same exponent, for every smooth positive probability form and Hermitian perturbation on $Y$.
\end{corollary}
\begin{proof}
Semipositivity gives nefness. If some $t_j$ is positive, compactness and strict positivity of the functions $f_j$ give
\[
 \begin{gathered}
 cB\leq\theta\leq CB,\qquad
 c=\min_{j:t_j>0}\bigl(t_j\min_X f_j\bigr)>0,\\
 C=\max_{j:t_j>0}\bigl(t_j\max_X f_j\bigr),
 \end{gathered}
\]
and \cref{prop:nef-maximum-rank-comparison} applies with $\varphi=0$. If all parameters vanish, then $B=\theta=0$, $R=0$, any $c,C>0$ may be used, and the order is $n$.

These inequalities between smooth forms are preserved by pullback, and the rank of $\pi^*B$ is at most $R$. If $R>0$, the set where $B$ has rank $R$ is nonempty and open, so it meets the dense open set where $\pi$ is biholomorphic; hence the maximal rank of $\pi^*B$ is $R$. This is immediate if $R=0$. Since $\dim Y=n$, the same proposition applies on $Y$, with any choice of data there. No singular potential needs to be pulled back.
\end{proof}

The constants depend on the parameters $t_j$ and need not stay bounded as a positive parameter tends to zero. The rank of a sum records how the generators overlap, and it need not equal the sum of their ranks. Members of the family are not closed when $\sum_jt_j\,\mathrm df_j\wedge B_j\ne0$.

\Needspace{10\baselineskip}
\begin{example}\label{ex:nef-open-zero-rank-background}
Let $X=\mathbb P^1\times\mathbb P^q$, $q\geq1$, with projections $p_1,p_2$. Choose a smooth nonnegative function $b\not\equiv0$ on $\mathbb P^1$ that vanishes on a nonempty open set, and a smooth strictly positive nonconstant function $f$ on $\mathbb P^q$. With $\omega_1$ the normalized Fubini--Study form on $\mathbb P^1$, set
\[
 B=p_1^*(b\omega_1),\qquad \theta=(f\circ p_2)B.
\]
The smooth nef background $\theta$ is not closed, vanishes on a nonempty open set, and has common order $q$, with two-sided bounds by multiples of $\epsilon^q$.
\end{example}
\begin{proof}
The form $b\omega_1$ is closed, because the first factor is a curve, and semipositive. Its pullback has rank one where $b>0$ and rank zero where $b=0$, so its maximal rank is one. Since $n=q+1$, \cref{cor:nef-variable-semipositive-generators} gives the order $q$. A nonconstant smooth function $f$ on the connected second factor has $\mathrm df\ne0$ somewhere. At a point of the product where $b>0$ and $\mathrm df\neq0$, the horizontal two-form $B$ wedges nontrivially with the vertical one-form $p_2^*\mathrm df$. Thus $\mathrm d\theta=p_2^*\mathrm df\wedge B\ne0$.
\end{proof}

These are geometric sufficient conditions for an integral common order, in the direction of \cref{qu:nef-degeneration-exponent}. They do not cover backgrounds without a closed semipositive comparison form, coefficient functions that are allowed to vanish, or the comparison of orders under residuals or non-K\"ahler modifications in general.
